\section{Abendmahl}
Wir feiern Weihnachten, weil um 200\nChr{} jemand die schöne Idee kam, den Geburtstag von Jesus Christus zu feiern.

Es gibt aber noch ein viel wichtigeres Fest, dass wir heute hier Feiern. Das Fest wurde nicht von Menschen erfunden und eingesetzt, sondern von Jesus Christus persönlich.

Dieses Fest, das Abendmahl, soll uns an den neuen Bund erinnern, den Gott mit uns und seinem Volk geschlossen hat.

Er hat es uns als ein Erinnerungsmahl eingesetzt. Ein Mahl, das uns daran erinnern soll, dass Jesus für unsere Sünden am Kreuz gestorben ist, dass er am 3. Tag von den Toten auferstanden ist und somit den Tod besiegt hat.

Ich lese zuerst die bekannten Worte aus
\biblezit{IKor}{10:23-28}

Wie bei uns üblich ist, warten wir, bis das Brot und der Wein ausgeteilt ist und nehmen es dann jeweils gemeinsam ein.

Würdest du bitte für das Brot beten.

\beten{Beten für das Brot}

{\textbf{Austeilen Brot}}

\textit{Wenn das Austeilen beendet ist mit diesen Worten, die Einnahme einleiten. \betonung[b, p]{Das Brot hochhalten!}}
\begin{bibelbox}{SCHL}{Lk}{22:19}
Und er nahm das Brot, dankte, brach es, gab es ihnen und sprach: Das ist mein Leib, der für euch gegeben wird; das tut zu meinem Gedächtnis!
\end{bibelbox}

\textit{Brot einnehmen und Gedenkminute}


Würdest du bitte für den Wein beten. 

\beten{Für den Wein}

{\textbf{Austeilen Wein}}

\textit{Wenn das Austeilen beendet ist mit diesen Worten, die Einnahme einleiten. \betonung[b, p]{Den Wein hochhalten!}}
\begin{bibelbox}{SCHL}{Lk}{22:20}
Desgleichen [nahm er] auch den Kelch nach dem Mahl und sprach: Dieser Kelch ist der neue Bund in meinem Blut, das für euch vergossen wird.
\end{bibelbox}

\textit{Wein einnehmen und Gedenkminute}